% =====================================================================
%  NetPredict — Complete Project Report & Interview Guide
%  Author: Ankit Kumar Patel
%  Compile: pdflatex (run twice) or drop into Overleaf (A4)
% =====================================================================
\documentclass[11pt,a4paper]{article}

\usepackage[T1]{fontenc}
\usepackage[utf8]{inputenc}
\usepackage{lmodern}
\usepackage[a4paper,margin=2.2cm]{geometry}
\usepackage{amsmath,amssymb}
\usepackage{booktabs}
\usepackage{tabularx}
\usepackage{array}
\usepackage{xcolor}
\usepackage{enumitem}
\usepackage{fancyhdr}
\usepackage{listings}
\usepackage{microtype}
\usepackage[hidelinks,colorlinks=true,linkcolor=linknavy,urlcolor=linknavy]{hyperref}
\Urlmuskip=0mu plus 1mu\relax
\makeatletter
\g@addto@macro\UrlBreaks{\do\-\do\_\do\/}
\makeatother

% ---- colours -------------------------------------------------------
\definecolor{linknavy}{HTML}{1A4E8A}
\definecolor{accent}{HTML}{1A4E8A}
\definecolor{soft}{HTML}{EEF3FA}
\definecolor{warnbg}{HTML}{FFF4E5}
\definecolor{codebg}{HTML}{F5F7FA}
\definecolor{greydk}{HTML}{444444}

% ---- headers / footers --------------------------------------------
\pagestyle{fancy}
\fancyhf{}
\lhead{\small\textcolor{greydk}{NetPredict — Project Report \& Interview Guide}}
\rhead{\small\textcolor{greydk}{Ankit Kumar Patel}}
\cfoot{\small\thepage}
\renewcommand{\headrulewidth}{0.3pt}

% ---- listings ------------------------------------------------------
\lstset{
  basicstyle=\ttfamily\footnotesize,
  backgroundcolor=\color{codebg},
  frame=single,
  rulecolor=\color{soft},
  breaklines=true,
  columns=fullflexible,
  keepspaces=true,
  showstringspaces=false,
  xleftmargin=4pt,
  framexleftmargin=2pt,
  aboveskip=6pt,belowskip=6pt
}

% ---- a simple "note" box (no exotic packages) ----------------------
\newcommand{\note}[1]{%
  \par\vspace{3pt}\noindent
  \colorbox{soft}{\parbox{\dimexpr\linewidth-2\fboxsep}{\small #1}}%
  \par\vspace{3pt}}
\newcommand{\warn}[1]{%
  \par\vspace{3pt}\noindent
  \colorbox{warnbg}{\parbox{\dimexpr\linewidth-2\fboxsep}{\small #1}}%
  \par\vspace{3pt}}

\setlist{itemsep=1pt,topsep=3pt}
\renewcommand{\arraystretch}{1.25}
\setlength{\emergencystretch}{3em}
\newcommand{\pf}[1]{\nolinkurl{#1}}

% small Q/A helpers
\newcommand{\Q}[1]{\vspace{4pt}\noindent\textbf{Q. #1}\par\nopagebreak}
\newcommand{\A}[1]{\noindent\textit{A.} #1\par}

\title{\vspace{-1.2cm}\textbf{\Large NetPredict}\\[2pt]
\large Complete Project Report \& Interview Preparation Guide\\[6pt]
\normalsize A predictive early-warning system for network congestion}
\author{Ankit Kumar Patel}
\date{November 2026}

\begin{document}
\maketitle
\thispagestyle{fancy}

\begin{abstract}
\noindent This document explains the \textbf{NetPredict} project end to end, in plain language.
It covers what the project is, \emph{why} it was chosen, how the data was \emph{generated},
how it was cleaned, which features were engineered, which models were trained and why,
how uncertainty and explainability were added, and what lives on \textbf{Hugging Face} and
\textbf{Kaggle}. It finishes with a glossary, a list of likely \textbf{cross-questions} a data-analytics
reviewer may ask (with answers), and an honest \textbf{``be careful about this''} section so there are
no surprises in a viva or interview. Read it once fully; then revise the glossary and the Q\&A.
\end{abstract}

\tableofcontents
\vspace{6pt}
\hrule

% =====================================================================
\section{One-paragraph summary (say this first)}
% =====================================================================
\noindent\emph{``NetPredict is a predictive early-warning system for computer networks. Normal monitoring
tools (Datadog, Zabbix, PRTG) are \textbf{reactive} --- they alert you \emph{after} packets have already
dropped. NetPredict is \textbf{proactive}: it reads live network signals (queue fill, latency, jitter,
loss), measures how fast they are rising, and forecasts the risk of congestion \textbf{5, 15 and 30 minutes
in advance}. On top of the forecast it adds three things that make it trustworthy: calibrated probabilities,
mathematically guaranteed uncertainty intervals (conformal prediction), and per-alert root-cause
explanations (SHAP). It is a full-stack app --- a Python/FastAPI backend, a React ``cockpit'' dashboard,
a documentation portal, plus published models, datasets and a live demo on Hugging Face.''}

\note{\textbf{30-second version for a non-technical examiner:} ``Think of a highway. Old monitoring tells
you after a 10-car pile-up. NetPredict notices the cars bunching up and warns you 15 minutes before the jam,
so engineers can reroute traffic before users feel anything.''}

% =====================================================================
\section{Why I chose this project (the motivation story)}
% =====================================================================
This is the ``why'' an examiner always asks. The honest, memorable version from the project's own dev diary:

\begin{itemize}
  \item \textbf{Personal pain point.} On the campus/hostel LAN, whenever someone pulled a large dataset or
  torrented, everyone else's SSH froze, voice calls turned robotic, and game ping jumped from 14\,ms to
  1800\,ms. This is a classic problem called \textbf{bufferbloat}.
  \item \textbf{Monitoring was useless.} Traditional SNMP polling runs every 30--60 seconds. By the time it
  said ``utilisation 98\%'', the queue had \emph{already} overflowed and dropped packets. The alarm arrived
  \emph{after} the damage.
  \item \textbf{The driving question.} ``Why detect a disaster \emph{after} it happens? Why not \emph{predict}
  the crash 5, 15 or 30 minutes early?'' --- that question became NetPredict.
  \item \textbf{Wanted something real, not another CRUD-over-an-LLM project.} The goal was to touch
  ``bare metal and mathematical reality'' --- queuing theory, kernel telemetry, time-series ML ---
  the kind of problem hyperscaler SRE teams (Cloudflare, Meta, Google) actually face.
\end{itemize}

\note{\textbf{Why it is a good data-analytics project:} it is a full analytics lifecycle --- messy raw data
$\rightarrow$ cleaning $\rightarrow$ feature engineering $\rightarrow$ honest validation $\rightarrow$
modelling $\rightarrow$ uncertainty $\rightarrow$ explainability $\rightarrow$ a decision/``what-if'' tool
$\rightarrow$ deployment. That is exactly the story a reviewer wants to see.}

% =====================================================================
\section{The problem domain in simple words}
% =====================================================================
\subsection{What is bufferbloat?}
Routers and switches have memory buffers (queues) that hold packets waiting their turn to be sent.
If buffers are too big and traffic exceeds the link's capacity, packets are not dropped immediately ---
they \emph{sit} in the queue for a long time. This drives round-trip time (RTT, i.e.\ latency) from a
healthy $\sim$15\,ms up to hundreds or thousands of ms, \emph{before} any packet is actually lost.
High latency with no loss yet = bufferbloat.

\subsection{The key insight: the delay gradient $\tfrac{dRTT}{dt}$}
The project's core idea is that \textbf{the rate of change of latency is a leading indicator}.
Using simple fluid/queuing physics: if the arrival rate $\lambda(t)$ exceeds link capacity $C$, the queue
grows at rate $\tfrac{dQ}{dt}=\lambda(t)-C$, and latency rises as
\[
\frac{dRTT}{dt}=\frac{1}{C}\big(\lambda(t)-C\big)=\rho(t)-1,
\qquad \rho=\frac{\lambda}{\mu}\ \text{(traffic intensity / utilisation)}.
\]
So a \emph{positive} slope means the queue is filling and a tail-drop ``cliff'' is coming. NetPredict watches
this slope, not just the raw level. In the diary's words: \emph{``the gradient is the leading indicator;
packet drop is merely the funeral notice.''}

\subsection{Reactive vs predictive (the headline contrast)}
\begin{center}
\begin{tabularx}{\linewidth}{@{}p{4.2cm}X@{}}
\toprule
\textbf{Traditional monitoring} & Alerts \emph{after} packets drop / buffer overflows (lagging). \\
\textbf{NetPredict} & Forecasts risk 5/15/30 min \emph{ahead} from trends, so engineers can reroute, rate-limit, or add buffer first. \\
\bottomrule
\end{tabularx}
\end{center}

% =====================================================================
\section{System architecture (the big picture)}
% =====================================================================
NetPredict is a full-stack system, not just a notebook.

\begin{center}
\begin{tabularx}{\linewidth}{@{}p{3.6cm}X@{}}
\toprule
\textbf{Layer} & \textbf{What it is} \\
\midrule
Backend & Python 3.12 + \textbf{FastAPI}. ``Hexagonal'' (ports \& adapters) design: data collectors, a sliding-window state buffer, feature extraction, ML models, and REST API routers are cleanly separated. \\
Frontend & \textbf{React 19 + TypeScript + Vite} ``cockpit'' dashboard (7 sections). Also has a client-side fallback engine so it runs free in a browser with no backend. \\
Docs portal & \textbf{Next.js} documentation site (Fumadocs) + an Obsidian knowledge vault. \\
Packaging & \textbf{Docker Compose} to run backend + frontend together. \\
Published & Models + datasets + a live Gradio demo on \textbf{Hugging Face}; a \textbf{Kaggle} notebook. \\
\bottomrule
\end{tabularx}
\end{center}

\subsection{Data plane (how a reading flows through the system)}
\begin{enumerate}
  \item \textbf{Ingest} a telemetry record (bandwidth, queue, RTT, jitter, loss, \dots) from one of three
  sources: a Linux \textbf{eBPF} kernel probe, a cross-platform \textbf{psutil} collector (reads your real
  NIC counters + measures live socket RTT to 1.1.1.1:53 / 8.8.8.8:53), or a \textbf{synthetic} trace.
  \item \textbf{Sliding window buffer:} a thread-safe ring buffer, \texttt{collections.deque(maxlen=720)} per
  interface (= 12 hours of 1-minute ticks). $O(1)$ append/evict, bounded memory, no GC stalls.
  \item \textbf{Causal feature extractor:} turns the window into rolling/slope/EWMA features
  (strictly backward-looking --- no peeking at the future).
  \item \textbf{Two decoupled ML engines run in parallel:}
  Stage 1 = Isolation Forest (``what is abnormal \emph{now}?''),
  Stage 2 = multi-horizon LightGBM (``what fails \emph{next}, at T+5/15/30?'').
  \item \textbf{Trust layer:} TreeSHAP explains each alert; conformal predictor adds a 90\% uncertainty band.
  \item \textbf{Action layer:} a counterfactual ``what-if'' simulator lets an operator test a fix (reroute/rate-limit/bigger buffer) \emph{before} touching real switches.
\end{enumerate}

% =====================================================================
\section{The dataset --- how it was generated (most important section)}
% =====================================================================
\warn{\textbf{Be clear and honest about this in an interview.} The training data is \textbf{synthetic}
(computer-generated with a physics-grounded simulator), \emph{not} captured from a real production network.
This is a \emph{strength}, not a weakness, if you explain it well: synthetic data gives you exact ground-truth
labels, full control over rare failure events, and perfect reproducibility (fixed random seed). Say it
plainly --- do not claim it is real captured traffic.}

\subsection{What the generator does (GitHub backend dataset)}
File: \texttt{backend/app/services/dataset\_loader.py}, function
\texttt{generate\_high\_fidelity\_network\_trace()}. It builds a \textbf{7-day, 1-minute} trace =
$7\times1440 = \mathbf{10{,}080}$ rows, with \texttt{np.random.seed(42)} (reproducible), for one
\texttt{core-router-alpha}, interface \texttt{xe-0/0/1}, starting 2026-09-01. Step by step:

\begin{enumerate}
  \item \textbf{Daily (diurnal) traffic cycle.} A cosine wave gives a realistic day/night pattern ---
  trough at 04:00, peak at 16:00:
  \[ \text{diurnal}=0.5\big(1-\cos(2\pi(h-4)/24)\big),\quad
     \text{util}=25+45\cdot\text{diurnal}+\mathcal{N}(0,3.5),\ \text{clipped }[10,95]\%. \]
  \item \textbf{Throughput} tied to utilisation on a 10\,Gbps link:
  $\text{throughput}=\tfrac{\text{util}}{100}\times10000+\text{noise}$.
  \item \textbf{Queue occupancy} (bufferbloat core): builds once utilisation passes $\sim$40\%:
  $q=\max(0,(u-40)\times1.2)+\text{Exp}(1.5)$.
  \item \textbf{Latency from queue (quadratic blow-up):}
  $RTT=14+(q/100)^2\times120+\text{noise}$, floored at 11\,ms. This reproduces the bufferbloat curve.
  \item \textbf{Packet loss / discards / retransmits} switch on sharply only when the queue passes 85\%:
  $\text{loss}=((q-85)/15)^{1.8}\times8.5+\text{Exp}(0.4)$; otherwise near zero. This is the ``cliff''.
  \item \textbf{Nine scripted incidents} are injected over the week (microburst surge, bufferbloat collapse,
  BGP reroute storm, interface CRC drop, memory-leak exhaustion) so the data contains realistic, labelled failures.
  \item \textbf{CPU / memory / BGP-flap} signals are derived from utilisation, queue and incidents.
  \item \textbf{Ground-truth labels (the targets):} an event is ``severe'' if
  (queue$>$80\% OR loss$>$1.5\% OR RTT$>$65\,ms). Then the \emph{future} labels are made by looking ahead:
  \texttt{target\_t5/t15/t30} = ``is it severe 5/15/30 minutes from now?'' and
  \texttt{future\_rtt\_t15}, \texttt{future\_loss\_t15} are the actual future values (via a time shift).
\end{enumerate}

\note{\textbf{``MAWI/CAIDA calibrated'' --- what that honestly means.} The generator was \emph{designed to
resemble} the statistical shape of well-known public traffic traces (MAWI, CAIDA) --- diurnal cycles, heavy-tailed
bursts, queue physics. It does \emph{not} load real MAWI/CAIDA files. If asked, say: ``It's a physics- and
pattern-calibrated \emph{simulator} inspired by MAWI/CAIDA characteristics.''}

\subsection{Schema of the GitHub dataset (12 signals + labels)}
\begin{center}
\begin{tabularx}{\linewidth}{@{}p{3.6cm}p{2.1cm}X@{}}
\toprule
\textbf{Column} & \textbf{Type/Unit} & \textbf{Meaning} \\
\midrule
\texttt{timestamp} & datetime & 1-minute grid \\
\texttt{device\_id},\ \texttt{interface\_id} & string & which router/port \\
\texttt{bandwidth\_util\_pct} & \% & link load \\
\texttt{throughput\_mbps} & Mbps & actual data rate \\
\texttt{rtt\_ms} & ms & latency (round-trip time) \\
\texttt{rtt\_jitter\_ms} & ms & latency stability \\
\texttt{queue\_occupancy\_pct} & \% & buffer fill (bufferbloat) \\
\texttt{packet\_loss\_pct} & \% & dropped packets \\
\texttt{tcp\_retrans\_rate} & /s & TCP strain \\
\texttt{interface\_discards\_sec} & /s & hardware tail-drops \\
\texttt{crc\_errors\_sec} & /s & physical/fibre errors \\
\texttt{cpu\_util\_pct},\ \texttt{memory\_util\_pct} & \% & host saturation \\
\texttt{bgp\_flap\_count} & int & routing instability \\
\texttt{target\_t5/t15/t30} & 0/1 & \textbf{labels}: severe in 5/15/30 min \\
\texttt{future\_rtt\_t15},\ \texttt{future\_loss\_t15} & ms, \% & \textbf{labels}: future values \\
\bottomrule
\end{tabularx}
\end{center}

\warn{\textbf{Class imbalance (expect this question).} Only about \textbf{2.8\%} of rows are positive for
\texttt{target\_t15}. Failures are \emph{rare}. That is why plain accuracy is misleading and why we use
PR-AUC, ROC-AUC, Brier score, and \texttt{class\_weight="balanced"}.}

% =====================================================================
\section{Data cleaning pipeline (the analytics ``before/after'')}
% =====================================================================
The notebook starts from a deliberately \textbf{dirty} file,
\texttt{raw\_telemetry\_trace.csv} ($\approx$10{,}122 rows), which has real-world style defects injected on
purpose so the cleaning step is meaningful. Counts measured in the raw file:

\begin{center}
\begin{tabularx}{\linewidth}{@{}Xrp{6.4cm}@{}}
\toprule
\textbf{Defect (before)} & \textbf{Count} & \textbf{Fix applied} \\
\midrule
Duplicate timestamps & 42 & sort by time, drop duplicates \\
Missing values (NaN) & 179 & forward-fill then back-fill (\texttt{ffill/bfill}) \\
Negative RTT error codes ($-999$) & 17 & set to NaN, then interpolate \\
Buffer overflow spikes ($>$100\%) & 15 & set to NaN (32-bit counter wrap-around), clamp 0--100 \\
\bottomrule
\end{tabularx}
\end{center}

\textbf{Result:} a clean \texttt{telemetry\_trace.csv} with exactly \textbf{10{,}080} rows, no NaNs, no
duplicates, all percentages in $[0,100]$, RTT $\geq 1$\,ms. The cleaning logic (in order): chronological sort
(fixes out-of-order UDP arrivals) $\rightarrow$ de-duplicate $\rightarrow$ replace sensor error codes with NaN
$\rightarrow$ time-series imputation $\rightarrow$ physical-range clamping.

\note{\textbf{Why ffill/bfill (not mean/median)?} Network state is continuous and autocorrelated in time ---
the router's condition one minute ago is the best estimate of a missing reading now. Mean/median imputation
would erase the time structure. Forward-fill respects causality (uses the past), and a final back-fill only
patches a possible gap at the very start.}

% =====================================================================
\section{Feature engineering (turning signals into predictors)}
% =====================================================================
Raw gauges are not enough; the model needs \emph{trend} and \emph{interaction} features. The notebook adds
\textbf{17} engineered columns (the production backend extractor computes a richer set; see the caveat on
counts below). Grouped:

\begin{itemize}
  \item \textbf{Rolling windows} (short 5-min, medium 15-min): means, standard deviations, maxima, sums of
  RTT, queue, utilisation, loss, discards. These capture microbursts and sustained pressure.
  \item \textbf{Rate-of-change / slopes} (the delay-gradient idea):
  \texttt{rtt\_slope\_5m}, \texttt{queue\_slope\_5m}, \texttt{util\_slope\_5m} = \texttt{diff(5)/5}.
  \item \textbf{EWMA} (exponentially weighted moving average, $\alpha=0.2$): smooths jitter while still
  reacting to real trends --- \texttt{rtt\_ewma}, \texttt{queue\_ewma}.
  \item \textbf{Domain interaction features:}
  \[ \text{buffer\_stress\_index}=\frac{\text{queue}\times\text{util}}{100}, \]
  \[ \text{congestion\_pressure\_score}=\mathrm{clip}\big(0.4\,Q+5\max(0,\text{rtt\_slope})+10\,L+0.8\,D,\ 0,\ 150\big). \]
\end{itemize}

\warn{\textbf{Zero future leakage.} Every rolling/slope/EWMA feature looks \emph{only backward} in time
(\texttt{min\_periods=1}, past differences). It never uses values from after the current minute. This is
critical for an honest time-series model (see next section).}

\note{\textbf{Feature-count caveat (a reviewer may notice).} Different documents say 17 / 18 / 29 / 48
features. The reason: the Jupyter notebook adds \textbf{17} engineered columns on top of the raw signals; the
production backend's \texttt{feature\_extractor.py} computes a larger set (docs quote ``29''); an early diary
note said 48. If asked, give the concrete number for the artefact you are showing (``the notebook uses 17
engineered features on top of the cleaned signals'') rather than a vague figure.}

% =====================================================================
\section{Validation: walk-forward split with embargo (the honesty engine)}
% =====================================================================
\subsection{The mistake that teaches the lesson}
The first model scored \textbf{ROC-AUC 0.994} --- and it was \emph{wrong}. It used
\texttt{KFold(shuffle=True)}. In time series, neighbouring minutes are $\sim$99\% correlated, so a shuffled
split lets the model ``peek'' at samples next to the test point. The model was \emph{memorising}, not
forecasting. This is \textbf{temporal data leakage}.

\subsection{The fix: walk-forward + embargo gap}
Split the timeline \emph{chronologically} (never shuffle), and leave a \textbf{30-minute embargo gap} between
folds so that rolling features (e.g.\ a 15-min window) cannot straddle the boundary:
\begin{center}
\begin{tabularx}{\linewidth}{@{}p{3.3cm}p{3.1cm}X@{}}
\toprule
\textbf{Fold} & \textbf{Rows} & \textbf{Purpose} \\
\midrule
Train (65\%) & 0--6{,}500 & fit the model \\
Embargo & 30 min gap & block leakage \\
Calibration (15\%) & 6{,}530--8{,}000 & calibrate probabilities + conformal quantile \\
Embargo & 30 min gap & block leakage \\
Test (20\%) & 8{,}030--10{,}080 & honest out-of-sample score \\
\bottomrule
\end{tabularx}
\end{center}
After this fix the score dropped to a \emph{truthful} ROC-AUC of about \textbf{0.78--0.84}. Lower, but real.

\note{\textbf{The one-liner to remember:} ``Shuffling a time series inflates your score because the future
leaks into the past. Walk-forward validation with an embargo gap is the only honest way to test a forecaster.''}

% =====================================================================
\section{The models and why each was chosen}
% =====================================================================
\subsection{Why LightGBM (gradient-boosted trees), not an LSTM}
LSTMs/GRUs/Transformers were tried first and lost, for three concrete reasons:
\begin{enumerate}
  \item \textbf{Speed:} LSTM inference was $\sim$18\,ms; at 10{,}000 events/s that creates a backlog.
  LightGBM is sub-millisecond on a CPU, no GPU needed.
  \item \textbf{Tabular friction:} telemetry mixes continuous gauges with discrete counters/IDs. Trees handle
  mixed tabular features natively, with no scaling and graceful handling of missing values.
  \item \textbf{Sharp transients:} LSTM gating smooths over the sudden microburst spikes that matter most;
  trees cut sharp thresholds and keep the cliff edges.
\end{enumerate}
\textbf{Key reframe:} with a rich rolling window, the network state is approximately \textbf{Markovian} ---
you don't need recurrence; you can hand explicit trend features to a tree model.

\subsection{Primary model: Multi-Horizon LightGBM classifier}
One classifier per horizon (T+5, T+15, T+30). Hyperparameters:
\begin{lstlisting}
lgb.LGBMClassifier(n_estimators=120, max_depth=5, learning_rate=0.06,
                   num_leaves=24, class_weight="balanced",
                   random_state=42, verbose=-1)
\end{lstlisting}

\subsection{Baselines (to prove the model earns its place)}
\begin{itemize}
  \item \textbf{Heuristic:} alert if \texttt{queue\_occupancy > 75\%}. Simple rule; many false alarms.
  \item \textbf{Logistic Regression} (scaled features): linear, can't capture the non-linear queue curve.
  \item \textbf{Naive persistence / moving average:} ``it'll be like now / like the 15-min average''.
\end{itemize}

\subsection{Stage-1 anomaly detector (decoupled): Isolation Forest}
Unsupervised ``is this reading weird \emph{right now}?'' (\texttt{n\_estimators=100},
\texttt{contamination=0.05}, with a RobustScaler). It is deliberately \emph{separate} from the forecaster so
the dashboard can tell apart a harmless transient spike (anomaly, but low future risk) from a real impending
failure (high future risk). This 4-quadrant view kills alert fatigue.

\subsection{Continuous forecaster: LightGBM regressor}
Predicts the actual future RTT at T+15 (\texttt{n\_estimators=100, max\_depth=5, lr=0.07}); its residuals feed
the conformal interval.

\subsection{Second research track: TraceGNN (on Hugging Face)}
A separate, more advanced model: a directed \textbf{Graph Attention Network (GAT, 4 heads)} that takes a
microservice distributed-trace DAG (from OpenTelemetry) and finds the \emph{root-cause} service of a failure.
Reported top-1 accuracy 98.4\%, F1 0.978, ROC-AUC 0.994 vs weaker baselines (max-duration heuristic 64.2\%,
random forest 78.5\%). This shows graph ML in addition to tabular ML.

% =====================================================================
\section{Probability calibration (making the \% trustworthy)}
% =====================================================================
Raw tree outputs are good \emph{rankings} but poor \emph{probabilities} --- a ``0.85'' might really fail only
40\% of the time. If operators can't trust the number, they ignore alarms (alert fatigue). Fix:
\textbf{Isotonic Regression} (a monotonic, non-parametric fit solved by the Pool-Adjacent-Violators
Algorithm, PAVA) trained \emph{on the calibration fold}. Quality is measured by the \textbf{Brier score}
(mean squared error of probabilities, lower is better):
\[ BS=\frac{1}{N}\sum_{i=1}^{N}(f_i-y_i)^2. \]
After isotonic calibration, a predicted ``80\%'' really means $\sim$80 of 100 such cases fail.

\warn{\textbf{Numbers differ across the project's own docs} --- know this before the viva. The technical report
quotes Brier \textbf{0.0482 $\rightarrow$ 0.0157} (raw $\rightarrow$ calibrated); the dev diary quotes
\textbf{0.284 $\rightarrow$ 0.061}. The \emph{method} is identical; only the quoted run differs. Safest answer:
``Isotonic calibration cut the Brier score by roughly 3--4$\times$; the headline figure I report is
\textbf{0.0157}.''}

% =====================================================================
\section{Conformal prediction (guaranteed uncertainty bands)}
% =====================================================================
A single number (``RTT will be 142\,ms'') is reckless during a surge. \textbf{Split conformal prediction}
gives an interval with a \emph{distribution-free, finite-sample} coverage guarantee.

\textbf{Recipe (very simple):}
\begin{enumerate}
  \item On the calibration set, compute residuals $R_i=|y_i-\hat{y}_i|$ (how wrong the model was).
  \item Pick the rank $k=\lceil (n+1)(1-\alpha)\rceil$ and take $\hat q = R_{(k)}$ (the $(1-\alpha)$ empirical
  quantile). For 90\% coverage, $\alpha=0.10$.
  \item For any new point: interval $=[\hat y-\hat q,\ \hat y+\hat q]$.
\end{enumerate}
\textbf{Guarantee:} $P\big(y\in[\hat y-\hat q,\hat y+\hat q]\big)\geq 1-\alpha$, needing only that
calibration and test points are \emph{exchangeable} (no Gaussian assumption). In the backend
$\hat q_{RTT}\approx\pm6.5$\,ms and $\hat q_{loss}\approx\pm0.4\%$; measured test coverage $\approx$90\%.

\subsection{Three conformal ``flavours'' (shown in the Hugging Face demo)}
\begin{center}
\begin{tabularx}{\linewidth}{@{}p{3.5cm}Xp{2.4cm}@{}}
\toprule
\textbf{Method} & \textbf{What it guarantees} & \textbf{Drift-robust?} \\
\midrule
Split Conformal (Vovk 2005) & marginal coverage $P(y\in C(X))\geq1-\alpha$ & no \\
Adaptive CI --- ACI (Gibbs \& Cand\`es 2021) & long-run coverage under drift, via $\alpha_{t+1}=\alpha_t+\gamma(\alpha-\text{err}_t)$ & yes \\
Mondrian CP (Vovk 2003) & \emph{conditional} coverage \emph{per regime} (linear / microburst / saturation) using a separate $\hat q_k$ per group & partition-aware \\
\bottomrule
\end{tabularx}
\end{center}
\textbf{Why Mondrian here:} plain split-CP \emph{under-covers} during saturation (the dangerous regime) and
\emph{over-covers} in calm periods. Splitting the residuals by regime restores honest coverage exactly where
safety matters most --- during heavy-tailed microbursts.

% =====================================================================
\section{Explainability: TreeSHAP}
% =====================================================================
Operators won't trust a black box. \textbf{SHAP} (SHapley Additive exPlanations) splits a prediction into
exact additive contributions of each feature ($\phi_i$). \textbf{TreeSHAP} computes these in polynomial time
$O(T\cdot L\cdot D^2)$ --- fast enough for real time ($<$2\,ms) --- instead of KernelSHAP's exponential cost.
The dashboard renders a waterfall per alert, e.g.\ \emph{``queue-delay slope contributed +42\% to this
risk''}, so an engineer sees the root cause in milliseconds. Globally, mean $|$SHAP$|$ ranks the top drivers
of failure.

% =====================================================================
\section{What-if (counterfactual) simulation --- the decision tool}
% =====================================================================
A forecast is only useful if you can act. The \textbf{What-If Scenario Lab} lets an operator test a fix
\emph{before} touching production:
\begin{itemize}
  \item Sliders: \textbf{traffic diversion \%}, \textbf{rate-limit \%}, \textbf{buffer-expansion factor}.
  \item It recomputes the physical features (lower utilisation $\rightarrow$ lower queue $\rightarrow$ gentler
  slope), re-runs the calibrated model, and reports the new risk and a verdict
  (\textsc{optimal} $\leq-50\%$, \textsc{effective} $\leq-20\%$, \textsc{marginal}, \textsc{ineffective}).
\end{itemize}
Example from the notebook: an incident at risk \textbf{92.5\%}; simulating a \textbf{25\% traffic diversion}
drops the predicted risk into a safe state --- all \emph{before} reconfiguring any switch.

% =====================================================================
\section{Results \& benchmarks (with honest caveats)}
% =====================================================================
Out-of-sample test (T+15 horizon, $\sim$2{,}016 rows, incident rate $\approx$2.5\%):
\begin{center}
\begin{tabularx}{\linewidth}{@{}Xcccc@{}}
\toprule
\textbf{Model} & \textbf{PR-AUC} & \textbf{ROC-AUC} & \textbf{Brier} & \textbf{F1} \\
\midrule
LightGBM (calibrated) & 0.6125 & 0.8389 & \textbf{0.0157} & 0.706 \\
Logistic Regression   & 0.6157 & 0.8393 & 0.0166 & 0.649 \\
Moving Average (15m)  & 0.4866 & 0.8783 & 0.0724 & 0.413 \\
Naive Persistence     & 0.7096 & 0.8491 & 0.0329 & 0.706 \\
\bottomrule
\end{tabularx}
\end{center}

\warn{\textbf{Don't oversell --- and be ready for this.} On this particular test slice some baselines are
\emph{close to} (or on one metric better than) LightGBM. The honest framing: ``LightGBM's advantage is the
\textbf{calibrated, well-behaved probabilities} (best Brier by a wide margin) plus explainability and
uncertainty --- not a huge raw-AUC gap on an easy, highly imbalanced slice. On harder regimes and across
horizons the tree model generalises better than the linear/heuristic baselines.''}

\subsection{Latency (speed) numbers}
\begin{center}
\begin{tabularx}{\linewidth}{@{}Xcc@{}}
\toprule
\textbf{Component} & \textbf{p50} & \textbf{p99} \\
\midrule
Raw LightGBM C-booster (\texttt{predict}) & 0.048\,ms & 0.18\,ms \\
LightGBM via pandas path & 1.83\,ms & 2.94\,ms \\
TreeSHAP explainer & 2.95\,ms & 3.99\,ms \\
Conformal interval & 0.003\,ms & 0.008\,ms \\
End-to-end telemetry tick & 93\,ms & 119\,ms \\
\bottomrule
\end{tabularx}
\end{center}
\note{\textbf{Reconcile the ``0.14\,ms'' claim.} The sub-0.2\,ms figure is the \emph{native C-booster} path on
raw NumPy arrays. The $\sim$1.8\,ms is the convenience pandas path (DataFrame overhead). The 93\,ms end-to-end
tick includes everything else (feature build, anomaly scoring, SHAP, I/O). All three are true --- just be
precise about \emph{which} you quote.}

\subsection{The Hugging Face queuing-delay model's metrics}
The published model card for \texttt{netpredict-telemetry-lgbm} reports (on the 6-feature eBPF dataset):
RMSE 0.398\,ms, MAE 0.281\,ms, $R^2=0.9924$, 95\% conformal coverage 95.12\%, p50 inference 0.14\,ms,
$\hat q=0.742$\,ms. These describe a \emph{different} task (single-horizon queuing-delay \emph{regression}),
not the multi-horizon classifier above --- keep the two straight.

% =====================================================================
\section{Everything that lives on Hugging Face (your inventory)}
% =====================================================================
You asked specifically about this --- here is the full list under the account
\href{https://huggingface.co/The-AnkitPatel}{\texttt{The-AnkitPatel}}.

\begin{center}
\begin{tabularx}{\linewidth}{@{}p{4.9cm}p{2.1cm}X@{}}
\toprule
\textbf{Repo} & \textbf{Type} & \textbf{What it is} \\
\midrule
\pf{netpredict-telemetry-lgbm} & Model & LightGBM queuing-delay regressor (6 features $\rightarrow$ latency), MIT. \\
\pf{otel-trace-gnn} & Model & GAT graph model for microservice root-cause, Apache-2.0. \\
\pf{ebpf-network-telemetry} & Dataset & 15{,}000 rows, 6 eBPF/XDP features + latency target. \\
\pf{otel-microservice-trace-anomalies} & Dataset & 180{,}000 spans / 15{,}000 trace DAGs, 12 services, 5 fault types + nominal. \\
\pf{netpredict-conformal-telemetry} & Space (Gradio) & Live 7-tab demo of the whole conformal pipeline. \\
\bottomrule
\end{tabularx}
\end{center}

\subsection{The Hugging Face eBPF dataset --- how \emph{it} is generated}
This is a \emph{second, simpler} synthetic schema (6 features), generated in the Space's \texttt{app.py}
(\texttt{generate\_calibration\_data}, \texttt{np.random.default\_rng(42)}). Features:
\texttt{ingress\_mbps, active\_flows, arrival\_kpps, buffer\_util\_pct, loss\_rate\_pct, jitter\_ms};
target \texttt{latency\_ms}. The physics is an \textbf{M/M/1-style queuing} formula:
\[ \text{queuing}=\frac{\text{buffer\_util}}{100-\text{buffer\_util}+\varepsilon}\times0.85,\qquad
   \text{latency}=1.2+\text{queuing}+3.5\cdot\text{loss}+\text{noise}. \]
Loss turns on above 80\% buffer; each row is tagged with a \textbf{regime}
(linear / microburst / saturation) used by Mondrian CP. The live app synthesises 6{,}000 rows (60/40
train/calibration); the published CSV has 15{,}000 rows.

\subsection{The Space's 7 tabs (what the demo shows)}
(1) Split Conformal, (2) Adaptive CI, (3) Mondrian CP, (4) eBPF/XDP kernel pipeline (an XDP program parses
Ethernet/IP/TCP/UDP in a 1\,MiB ring buffer and exports the 6-feature schema in kernel space),
(5) \textbf{Treelite + clang -O3} C-compilation (drops the tree inference from $\sim$0.53\,ms Python to
\textbf{65--95\,ns} per call), (6) enterprise serving (Truss / Baseten), (7) preprint abstract + BibTeX +
benchmark table (mean 0.577\,ms, p50 0.532, p99 1.027\,ms, 1{,}731 RPS, sub-ms SLA 98.78\%).

\subsection{Also published}
A \textbf{Kaggle} notebook: \texttt{theankitkumarpatel/conformalized-lightgbm-network-telemetry}.

% =====================================================================
\section{Technology stack (one-glance table)}
% =====================================================================
\begin{center}
\begin{tabularx}{\linewidth}{@{}p{3.4cm}X@{}}
\toprule
\textbf{Area} & \textbf{Tools} \\
\midrule
Language & Python 3.12, TypeScript \\
Backend & FastAPI, Uvicorn, Pydantic \\
ML / stats & LightGBM, scikit-learn (IsolationForest, IsotonicRegression, LogisticRegression), SHAP, NumPy, pandas, SciPy \\
Data viz & Matplotlib, Seaborn, Plotly \\
Kernel telemetry & eBPF / XDP (BCC), psutil, raw sockets \\
Frontend & React 19, Vite, TypeScript \\
Docs & Next.js + Fumadocs, Obsidian vault \\
Serving / deploy & Docker Compose, Gradio (HF Space), Treelite + clang, Truss/Baseten, Vercel \\
Reproducibility & fixed seed 42, walk-forward split, versioned model artefacts (\texttt{.joblib}) \\
\bottomrule
\end{tabularx}
\end{center}

% =====================================================================
\section{Glossary (say these in plain words)}
% =====================================================================
\begin{description}[leftmargin=0pt,style=nextline]
  \item[Bufferbloat] Oversized buffers hold packets instead of dropping them, so latency balloons before any loss.
  \item[RTT / latency] Round-trip time --- how long a packet takes to go and come back.
  \item[Jitter] Variation in latency; unstable delay.
  \item[Delay gradient $dRTT/dt$] How fast latency is rising --- the early-warning signal.
  \item[Queue occupancy] How full the router buffer is (\%).
  \item[Diurnal cycle] The daily up-and-down pattern of traffic (quiet at night, busy afternoon).
  \item[Multi-horizon] Predicting several future times at once (5, 15, 30 min).
  \item[LightGBM / GBDT] Gradient-boosted decision trees --- many small trees added in sequence.
  \item[Isolation Forest] Unsupervised outlier detector (isolates weird points fast).
  \item[Temporal leakage] Letting future information sneak into training; inflates scores dishonestly.
  \item[Walk-forward split] Train on past, test on future, in time order (never shuffle).
  \item[Embargo gap] A blackout gap between folds so rolling features can't straddle the boundary.
  \item[Calibration / Brier] Making predicted \% match real frequency; Brier = MSE of probabilities.
  \item[Isotonic regression / PAVA] A monotonic calibration curve fit by pool-adjacent-violators.
  \item[Conformal prediction] Distribution-free prediction \emph{intervals} with a coverage guarantee.
  \item[Coverage] Fraction of true values that land inside the interval (target 90\%).
  \item[Mondrian CP] Conformal with a separate quantile per group/regime.
  \item[ACI] Adaptive conformal that adjusts $\alpha$ over time to stay honest under drift.
  \item[SHAP / TreeSHAP] Exact additive feature-attribution for each prediction; TreeSHAP is the fast tree version.
  \item[Counterfactual] A ``what if I changed X'' simulation.
  \item[M/M/1/K queue] A finite-buffer queuing model (Poisson arrivals/service, buffer size $K$); explains tail-drop.
  \item[eBPF / XDP] Linux tech to run small programs in the kernel to capture packet/TCP metrics with almost no overhead.
  \item[OpenTelemetry / span / trace / DAG] Standard for distributed tracing; a trace is a graph of service calls.
  \item[GAT] Graph Attention Network --- a neural net that learns on graphs using attention over neighbours.
\end{description}

% =====================================================================
\section{Likely cross-questions (with answers)}
% =====================================================================
These are the questions a data-analytics faculty / interviewer is most likely to fire. Short, confident answers.

\subsection{Data \& analytics}
\Q{Is this real data or synthetic?}
\A{Synthetic, by design --- a physics-grounded simulator with fixed seed 42, so I get exact ground-truth
labels, control over rare failures, and full reproducibility. The shape mimics public MAWI/CAIDA traffic.}

\Q{How exactly did you generate it?}
\A{A diurnal cosine traffic cycle drives utilisation; queue builds past $\sim$40\% utilisation; latency is
$14+(q/100)^2\cdot120$; loss switches on sharply above 85\% queue; nine scripted incidents inject failures;
labels are the ``severe'' condition shifted 5/15/30 minutes into the future.}

\Q{Why is accuracy not your main metric?}
\A{Only $\sim$2.8\% of rows are positive --- a model that always says ``no failure'' gets $\sim$97\% accuracy
but is useless. So I use PR-AUC, ROC-AUC, Brier score and F1, and \texttt{class\_weight="balanced"}.}

\Q{How did you handle missing values and outliers?}
\A{Sort by time, drop duplicate timestamps, convert sensor error codes ($-999$ RTT, $>$100\% queue) to NaN,
then forward-fill + back-fill (because network state is continuous in time), then clamp to physical ranges.}

\Q{Why forward-fill instead of mean imputation?}
\A{Mean/median destroys the time structure. The most recent valid reading is the best estimate of the current
hidden state, and forward-fill only uses the past (causal).}

\subsection{Modelling}
\Q{Why LightGBM over deep learning (LSTM)?}
\A{Sub-ms CPU inference vs $\sim$18\,ms for LSTM; native handling of mixed tabular features and missing values;
and trees keep the sharp microburst transients that LSTM gates smooth away. With a rich rolling window the
problem is effectively Markovian, so recurrence isn't needed.}

\Q{Your first model got 0.994 AUC --- why did you not report that?}
\A{It was temporal leakage from a shuffled split. After switching to walk-forward validation with a 30-minute
embargo, the honest score was $\sim$0.78--0.84. I report the honest one.}

\Q{What is the embargo gap and why 30 minutes?}
\A{A blackout gap between train/calibration/test so that rolling features (up to 15-min windows) and the
forecast horizon can't overlap across the boundary; 30 min safely exceeds the longest window + horizon lag.}

\Q{What do your three baselines prove?}
\A{That a simple threshold rule and a linear model (and naive persistence) can't match the calibrated,
non-linear tree model --- especially on probability quality (Brier).}

\Q{Why a separate anomaly detector and a forecaster?}
\A{They answer different questions: ``abnormal now'' vs ``will fail soon''. Decoupling them lets me suppress
harmless transient spikes (anomaly but low future risk) and avoid alert fatigue.}

\subsection{Statistics / uncertainty}
\Q{What does conformal prediction guarantee, and under what assumption?}
\A{$P(y\in C(X))\geq 1-\alpha$ in finite samples, assuming calibration and test points are exchangeable. No
distributional (Gaussian) assumption is needed.}

\Q{Difference between calibration and conformal?}
\A{Calibration makes the \emph{probability} honest (80\% means 80\%). Conformal makes the \emph{interval}
honest (the true value is inside 90\% of the time). Different guarantees.}

\Q{Why Mondrian conformal instead of plain split conformal?}
\A{Plain split-CP only guarantees \emph{marginal} coverage; it under-covers in the dangerous saturation regime
and over-covers in calm periods. Mondrian uses a per-regime quantile to restore \emph{conditional} coverage
where safety matters.}

\Q{What is the Brier score and what was your improvement?}
\A{Mean squared error of predicted probabilities (lower = better). Isotonic calibration cut it by roughly
3--4$\times$ (headline 0.0157).}

\subsection{Systems / deployment}
\Q{How is it fast enough for real time?}
\A{LightGBM's native C booster predicts in $\sim$0.05\,ms (p99 $<$0.2\,ms); conformal adds microseconds;
and with Treelite + clang -O3 the tree compiles to $\sim$65--95\,ns per call. No GPU needed.}

\Q{What does the eBPF part do?}
\A{A small in-kernel XDP/BCC program reads TCP metrics (smoothed RTT, retransmits, loss) and packet counters
with almost no overhead, and exports the feature schema to user space via a ring buffer.}

\Q{Walk me through the architecture.}
\A{Ingest $\rightarrow$ $O(1)$ sliding-window ring buffer $\rightarrow$ causal feature extractor $\rightarrow$
(Isolation Forest + multi-horizon LightGBM) $\rightarrow$ isotonic calibration + conformal interval +
TreeSHAP $\rightarrow$ what-if simulator $\rightarrow$ React cockpit.}

\subsection{``Gotcha'' questions (be ready)}
\Q{The two datasets have different columns --- explain.}
\A{Yes: the backend/notebook dataset is a 12-signal router/SNMP schema (10{,}080 rows, multi-horizon
classification + regression); the Hugging Face \texttt{ebpf-network-telemetry} is a simpler 6-feature eBPF
schema (15{,}000 rows, single-horizon latency regression). Same idea, two framings for two deliverables.}

\Q{Your model card quotes $R^2=0.9924$ but your classifier AUC is only $\sim$0.84 --- contradiction?}
\A{No --- different tasks. $R^2=0.9924$ is the \emph{regression} of queuing delay on the eBPF dataset
(an easy, smooth target); 0.84 is the \emph{classification} of rare future failures on the router dataset
(hard, imbalanced). They measure different things on different data.}

\Q{Can you prove the Hugging Face model works?}
\A{Be honest: the model \emph{card} is published with metrics, but the weight file (\texttt{model.txt}) is not
yet uploaded to that repo --- the Space regenerates and trains the model at runtime. If asked to load it from
the Hub, that download will currently fail; the reproducible path is the Space/Kaggle code.}

% =====================================================================
\section{Be careful about this (honesty checklist before any viva)}
% =====================================================================
\warn{Read this list once so nothing catches you off guard. None of it is fatal --- it is about
\emph{how you phrase} things.}
\begin{enumerate}
  \item \textbf{Data is synthetic.} Never call it real captured traffic. Call it a physics-grounded simulator.
  \item \textbf{``IIT Roorkee / Distributed Systems \& Telemetry Lab'', ``NIT Uttarakhand'', ``CNCF
  OpenTelemetry-Go contributor''} appear on the published cards. Only claim affiliations/roles you can actually
  back up; otherwise describe it as a solo student research project (which the dev diary itself says it is).
  \item \textbf{Metric inconsistencies} exist across the docs (Brier 0.0482$\rightarrow$0.0157 vs
  0.284$\rightarrow$0.061; AUC 0.78 vs 0.84; features 17/18/29/48; ``9 incidents'' vs a comment saying ``18'';
  latency 0.14 vs 1.8 vs 93\,ms). Pick one consistent set of headline numbers and know \emph{which artefact}
  each belongs to.
  \item \textbf{HF model weight file missing.} \texttt{netpredict-telemetry-lgbm} currently has only the README
  (no \texttt{model.txt}); the live model is trained inside the Space. Say so if asked to load it.
  \item \textbf{Two different schemas / datasets} (12-feature router vs 6-feature eBPF) --- don't mix their
  numbers.
  \item \textbf{Repo name typo:} some clone instructions say \texttt{NetPreditct} (extra ``t''); the repo is
  \texttt{NetPredict}.
  \item \textbf{Baselines are competitive} on the easy test slice --- lead with calibration/uncertainty/
  explainability as the real differentiators, not a huge AUC gap.
\end{enumerate}

% =====================================================================
\section{How to run / reproduce (quick reference)}
% =====================================================================
\begin{lstlisting}
# Full app (Docker)
git clone https://github.com/The-AnkitPatel/NetPredict && cd NetPredict
docker compose up --build        # dashboard :5173, API :8000/docs

# Backend only
python -m venv .venv && source .venv/bin/activate
pip install -r backend/requirements.txt psutil
pytest backend/tests -v
uvicorn backend.app.main:app --port 8000 --reload

# The analytics notebook (the data-analytics deliverable)
notebooks/NetPredict_ML_Pipeline.ipynb   # 11 cells, end to end

# Hugging Face live demo
# https://huggingface.co/spaces/The-AnkitPatel/netpredict-conformal-telemetry
\end{lstlisting}

The notebook is the cleanest thing to present to a data-analytics reviewer: it runs top-to-bottom through
ingestion $\rightarrow$ cleaning (with a before/after audit table) $\rightarrow$ feature engineering
$\rightarrow$ EDA plots $\rightarrow$ walk-forward split $\rightarrow$ model vs baselines $\rightarrow$
calibration $\rightarrow$ conformal intervals $\rightarrow$ TreeSHAP $\rightarrow$ what-if simulation.

% =====================================================================
\section*{Appendix: file map}
\addcontentsline{toc}{section}{Appendix: file map}
% =====================================================================
\begin{center}
\begin{tabularx}{\linewidth}{@{}p{6.2cm}X@{}}
\toprule
\textbf{Path} & \textbf{Contents} \\
\midrule
\pf{backend/app/services/dataset_loader.py} & synthetic 7-day trace generator \\
\pf{backend/app/services/feature_extractor.py} & causal rolling / slope / EWMA features \\
\pf{backend/app/services/sliding_window.py} & $O(1)$ ring buffer (deque maxlen 720) \\
\pf{backend/app/ml/predictor_model.py} & multi-horizon LightGBM + calibrators \\
\pf{backend/app/ml/conformal_predictor.py} & split conformal intervals \\
\pf{backend/app/ml/anomaly_detector.py} & Isolation Forest \\
\pf{backend/app/ml/shap_explainer.py} & TreeSHAP attributions \\
\pf{backend/app/services/simulation_service.py} & counterfactual what-if engine \\
\pf{backend/app/collectors/ebpf_sockops.c} & eBPF/BCC kernel TCP probe \\
\pf{notebooks/NetPredict_ML_Pipeline.ipynb} & 11-stage analytics notebook \\
\pf{docs/NETPREDICT_TECHNICAL_RESEARCH_REPORT.md} & long technical report \\
\texttt{obsidian\_vault/05 .../Dev Diary.md} & the ``why/journey'' narrative \\
\bottomrule
\end{tabularx}
\end{center}

\vfill
\begin{center}\small\textcolor{greydk}{Prepared as an end-to-end study and interview guide for the NetPredict project.}\end{center}

\end{document}
